\begin{abstract}
La revisión manual de videos extensos dificulta localizar lenguaje dañino cuando los subtítulos contienen modismos, ironía, ruido y categorías poco frecuentes. Este trabajo diseña y evalúa, mediante ciencia del diseño, un moderador semiautomático para videos peruanos de YouTube: produce alertas ligadas a intervalos temporales y conserva la decisión final en una persona supervisora. El corpus v2.1 contiene 173\,240 chunks elegibles de 4\,906 videos, con cinco salidas aprendidas: \texttt{SEGURO} y cuatro daños multietiqueta; 14\,163 chunks presentan al menos un daño. La evaluación compara 28 modelos individuales y cinco ensembles sobre 10\,600 filas de validación y cinco pliegues agrupados por video. El criterio predeclarado usa como métrica primaria la exactitud balanceada (BA) de cualquier daño, protege el desempeño por categoría mediante macro-AUPRC y congela un promedio suave de regresión logística, cascada E5 y Qwen3--0.6B adaptado mediante LoRA. En validación, el ensemble obtiene BA 0,84, macro-AUPRC 0,55 y macro-F1 0,57; mejora al mejor individuo en 0,0086 de BA y 0,039 de macro-AUPRC, aunque el contraste pareado no establece superioridad estadística. En la apertura única del test con prevalencia natural (22\,684 chunks; 7,9\% con daño), la compuerta alcanza BA 0,85 y sensibilidad 0,91. La política selectiva deriva 35\% a revisión y, sobre el 65\% automático, obtiene BA 0,94. El sistema ofrece desempeño competitivo y trazable para priorización; la evidencia no autoriza bloqueo o sanción automáticos.
\end{abstract}

\begin{IEEEkeywords}
moderación semiautomática, procesamiento del lenguaje natural, clasificación multietiqueta, lenguaje dañino, YouTube, Qwen, LoRA, ensemble, clasificación selectiva
\end{IEEEkeywords}
